    % ==== Functional Security in AI Systems =============================================================================
    \chapter{Functional Security in AI Systems}
    \section{Introduction}
    \subsection{Functional Security}
    \subsection{Automotive Safety}
    \subsection{Relevant Standards (ISO 26262, IEC 61508, ISO 27001, EU directive 2001/95/EG, ISO 25119)}
    \section{Functional Safety Standard ISO 26262}
    \subsection{Introduction}
    \subsection{Automotive Safe Integrity Levels (ASIL)}
    \subsection{Recommended Techniques}
    \section{IT Security Standards}
    \subsection{ISO 27001}
    \subsection{ISO 15408}
    \subsection{ISO 21434}
    \subsection{SAE J3061}
    \subsection{AECQ}
    \section{Approaches}
    \subsection{Safe Failure Fraction (SFF)}
    \subsection{Diagnostic Coverage (DC)}
    \subsection{Hazard Analysis and Risk Assessment (HARA)}
    \subsection{Fault Tree Analysis (FTA)}
    \subsection{Failure Modes, Effects \& (Diagnostic, Criticality) Analysis (FME[C,D]A)}
    \section{Attacks and Defenses}
    \subsection{Cyber Attacks}
    \subsection{Physical Attacks}
    \subsection{MISRA C/C++ Guidelines}

    \section{Literature}
    \begin{itemize}
        \item Smith D.J., Simpson K. (2016). The Safety Critical Systems Handbook. (4th ed.) Elsevier
        \item Smith D.j. (2017). Reliability, Maintainability and Risk (9th ed.) Elsevier
        \item Rausand, M. (2014). Reliability of Safety-Critical Systems: Theory and Applications Wiley.
    \end{itemize}

